\section*{Bab \ref{ch:b3:topology} --- Topologi Umum}

\begin{solution}{exo:b3:topology:1}
(a) Ada empat \omterm{def:b3:topology:topology}{topologi} pada $\{a,b\}$: yang indiskret
$\{\varnothing, X\}$; yang diskret; dan dua \omterm{def:b3:topology:topology}{topologi} \emph{Sierpiński}
$\{\varnothing, \{a\}, X\}$ serta $\{\varnothing,
\{b\}, X\}$. Kedua yang terakhir \omterm{def:b3:topology:continuity}{homeomorfik} (tukarkan $a, b$): jadi ada tiga
kelas. Yang \omterm{def:b3:topology:connected}{terhubung}: semuanya kecuali yang diskret (sebab hanya
\omterm{def:b3:topology:topology}{topologi} diskret yang memuat \omterm{def:b3:topology:topology}{himpunan terbuka} sekaligus tertutup yang sejati dan tak kosong).
Yang \omterm{def:b3:topology:hausdorff}{Hausdorff}: hanya yang diskret (sebab pada yang lain, satu-satunya
\omterm{def:b3:topology:topology}{persekitaran} $b$ adalah $X$, atau demikian pula untuk $a$).

(b) Untuk $\Q$: \omterm{def:b3:topology:interior}{interiornya} $\varnothing$ (sebab setiap selang memuat
bilangan irasional), \omterm{def:b3:topology:interior}{penutupnya} $\R$ (karena \omterm{def:b3:topology:interior}{padat}), batasnya $\R$. Untuk $[0,1)
\cup \{2\}$: \omterm{def:b3:topology:interior}{interiornya} $(0,1)$, \omterm{def:b3:topology:interior}{penutupnya} $[0,1]\cup\{2\}$,
batasnya $\{0, 1, 2\}$. Untuk $\{1/n\}$: \omterm{def:b3:topology:interior}{interiornya} $\varnothing$,
\omterm{def:b3:topology:interior}{penutupnya} $\{1/n\}\cup\{0\}$, dan batasnya \omterm{def:b3:topology:interior}{penutup} itu sendiri.
\end{solution}

\begin{solution}{exo:b3:topology:2}
(a) Setiap $V \subseteq Y$ yang terbuka merupakan gabungan $\bigcup B_j$ himpunan
basisnya, dan $f^{-1}(V) = \bigcup f^{-1}(B_j)$: jadi jika yang belakangan
terbuka, maka yang terdahulu pun terbuka; konversnya trivial.

(b) Prapeta $f^{-1}\bigl(\intoo{\frac12}{\frac32}\bigr) = [0, \infty)$
tidak terbuka: jadi $f$ tidak \omterm{def:b3:topology:continuity}{kontinu}; sedangkan \omterm{def:b3:topology:continuity}{kekontinuan} kanan di setiap
titik jelas ($f$ konstan lokal di sebelah kanan setiap
titik). Himpunan $[a, b)$ memenuhi kriteria basisnya
(\cref{def:b3:topology:basis}): $[a,b)\cap[c,d) = [\max(a,c),
\min(b,d))$. \omterm{def:b3:topology:topology}{Topologi} Sorgenfrey lebih halus daripada \omterm{def:b3:topology:topology}{topologi}
biasanya, sebab $(a, b) = \bigcup_{a < c < b}[c, b)$; dan secara sejati:
$[0,1)$ terbuka-Sorgenfrey tetapi tidak terbuka-biasa. \omterm{def:b3:topology:continuity}{Kekontinuan} dari
garis Sorgenfrey di $x$: \omterm{def:b3:topology:topology}{persekitaran} Sorgenfrey dari $x$
memuat himpunan \omterm{def:b3:topology:basis}{basis} $[a, b) \ni x$, sehingga memuat $[x, b)$ ---
jadi syarat \omterm{def:b3:topology:continuity}{kekontinuannya} terbaca: untuk setiap $\varepsilon > 0$
ada $\delta > 0$ dengan $\abs{f(t) - f(x)} < \varepsilon$
bilamana $x \leq t < x + \delta$. Dan itu tepat merupakan
\omterm{def:b3:topology:continuity}{kekontinuan} kanan di $x$.
\end{solution}

\begin{solution}{exo:b3:topology:3}
Dua \omterm{def:b3:topology:topology}{himpunan terbuka} tak kosong berkomplemen berhingga, sehingga
irisannya berkomplemen berhingga: jadi tak kosong (karena himpunannya
tak berhingga) --- sehingga tak ada dua titik yang berpersekitaran saling lepas: jadi tidak
\omterm{def:b3:topology:hausdorff}{Hausdorff}, dan karenanya tidak dapat dimetrikan
(\cref{def:b3:topology:hausdorff}). Misalkan $(x_n)$ injektif dan
$x$ sembarang: \omterm{def:b3:topology:topology}{persekitaran} $x$ memuat \omterm{def:b3:topology:topology}{himpunan terbuka} kofinit
$U$; sedangkan berhingga banyak titik $X\setminus U$ dikenai paling
banyak berhingga indeks (karena injektif), sehingga sebuah ekor
barisannya terletak di $U$: jadi $x_n \to x$, untuk \emph{setiap} $x$. Bukti
ketunggalannya memerlukan dua \omterm{def:b3:topology:topology}{persekitaran} saling lepas untuk memisahkan
dua limit yang diklaim --- dan persis itulah yang gagal di sini.
\end{solution}

\begin{solution}{exo:b3:topology:4}
(a) \omterm{def:b3:topology:continuity}{Kekontinuannya} langsung menurut konstruksinya
(\cref{def:b3:topology:constructions}). Keterbukaannya: $W
\subseteq X \times Y$ yang terbuka merupakan gabungan kotak $U_i\times V_i$, dan
$\pi_X(W) = \bigcup U_i$ (sebab kotak tak kosong terproyeksikan pada
faktornya), yang terbuka. Tidak tertutup: $H = \{(x,y) : xy = 1\}$ tertutup
di $\R^2$ (sebagai prapeta $\{1\}$ di bawah perkalian yang \omterm{def:b3:topology:continuity}{kontinu}),
tetapi $\pi_X(H) = \R\setminus\{0\}$ tidak tertutup.

(b) ($\Rightarrow$) Proyeksinya \omterm{def:b3:topology:continuity}{kontinu}. ($\Leftarrow$)
Misalkan $x^{(k)} \to x$ komponen demi komponen dan $W = \prod U_n$ sebuah
\omterm{def:b3:topology:topology}{persekitaran} \omterm{def:b3:topology:basis}{basis} dari $x$, dengan $U_n = X_n$ kecuali untuk $n \in F$
yang berhingga. Untuk setiap $n \in F$, $x_n^{(k)} \in U_n$ bila $k \geq
K_n$; sehingga untuk $k \geq \max_F K_n$ berlaku $x^{(k)} \in W$. Rumus $d$
mendefinisikan sebuah metrik (sebab setiap sukunya metrik, sampai pada
pemeriksaan baku bahwa $\min(1, d_n)$ demikian); bolanya: $B_d(x,
2^{-N})$ memuat kotak \omterm{def:b3:topology:basis}{basis} $\{y : d_n(x_n,y_n) < 2^{-N},\
n \leq N'\}$ untuk $N'$ yang sesuai (sebab ekor $\sum_{n > N'} 2^{-n}$
kecil), dan sebaliknya setiap kotak \omterm{def:b3:topology:basis}{basis} memuat sebuah bola-$d$:
jadi kedua \omterm{def:b3:topology:topology}{topologi} itu berpersekitaran sama di setiap titik.

(c) Pada \omterm{def:b3:topology:constructions}{topologi hasil kali}, $x^{(k)} \to 0$ menurut (b). Namun himpunan
$U = \prod_n \intoo{-\tfrac1n}{\tfrac1n}$ yang terbuka-kotak memuat $0$, sedangkan
$x^{(k)} \notin U$ untuk setiap $k$ (sebab koordinat ke-$n$ bernilai $1/k
\geq 1/n$ bila $n \geq k$): jadi tak ada ekornya yang masuk $U$. \omterm{def:b3:topology:topology}{Topologi} kotak
lebih halus secara sejati dan bukan \omterm{def:b3:topology:topology}{topologi} bertipe hasil kali untuk
keperluan kekonvergenan.
\end{solution}

\begin{solution}{exo:b3:topology:5}
Pemetaan $\varphi \colon \R \to S^1$, $x \mapsto \eu^{2\iu\pi x}$, bersifat
\omterm{def:b3:topology:continuity}{kontinu}, surjektif, dan konstan pada kelas modulo $\Z$: sehingga ia
menginduksikan bijeksi \omterm{def:b3:topology:continuity}{kontinu} $\bar\varphi \colon \R/\Z \to
S^1$ (\cref{prop:b3:topology:universal}(b)). Proyeksi
$\pi$ bersifat \emph{terbuka}: sebab untuk $U \subseteq \R$ terbuka,
$\pi^{-1}(\pi(U)) = \bigcup_{n}(U + n)$ terbuka, sehingga $\pi(U)$
terbuka. Sifat \omterm{def:b3:topology:hausdorff}{Hausdorffnya}: misalkan $\bar x \neq \bar y$; pilihlah
wakil dengan $\delta = \min_{n}\abs{x - y - n} > 0$;
maka peta selang berjari-jari $\delta/3$ di sekitar $x$ dan
$y$ bersifat terbuka (karena $\pi$ terbuka), memuat $\bar x, \bar y$, dan
saling lepas (sebab dua titik prapetanya akan berjarak $< 2\delta/3$
modulo $\Z$). Kekompakannya: $\R/\Z = \pi([0,1])$, yaitu peta \omterm{def:b3:topology:continuity}{kontinu}
\omterm{def:b3:topology:compact}{ruang kompak} (\cref{thm:b3:topology:compactprops}(2)). Kini
$\bar\varphi$ merupakan bijeksi \omterm{def:b3:topology:continuity}{kontinu} dari \omterm{def:b3:topology:compact}{ruang kompak} ke
$S^1$ yang \omterm{def:b3:topology:hausdorff}{Hausdorff}: jadi sebuah \omterm{def:b3:topology:continuity}{homeomorfisma}
(\cref{cor:b3:topology:compacthomeo}).

Untuk (iii): komposisi $[0,1] \hookrightarrow \R
\xrightarrow{\pi} \R/\Z$ bersifat \omterm{def:b3:topology:continuity}{kontinu}, surjektif, dan
mengenali tepat $0 \sim 1$: sehingga ia menginduksikan bijeksi \omterm{def:b3:topology:continuity}{kontinu}
$[0,1]/(0\sim1) \to \R/\Z$ dari \omterm{def:b3:topology:compact}{ruang kompak} (yaitu peta
\omterm{def:b3:topology:continuity}{kontinu} $[0,1]$ di bawah proyeksi kuosiennya) ke ruang yang
\omterm{def:b3:topology:hausdorff}{Hausdorff}: jadi sebuah \omterm{def:b3:topology:continuity}{homeomorfisma}. Dengan mengomposisikannya: ketiga ruang itu
\omterm{def:b3:topology:continuity}{homeomorfik}.
\end{solution}

\begin{solution}{exo:b3:topology:6}
(a) Selimut $K_1 \cup \dots \cup K_m$ terbatas menjadi selimut
setiap $K_j$: jadi berhingga \omterm{def:b3:topology:topology}{himpunan terbuka} per kepingan sudah cukup. Sedangkan
irisan $\bigcap K_i$ tertutup di dalam $K_{i_0}$ yang \omterm{def:b3:topology:compact}{kompak}
(sebab himpunan \omterm{def:b3:topology:compact}{kompak} bersifat tertutup di ruang metrik lingkungannya), sehingga
\omterm{def:b3:topology:compact}{kompak}.

(b) Fungsi $x \mapsto d(x, F)$ bersifat \omterm{def:b3:topology:continuity}{kontinu} (Lipschitz-$1$) dan
positif sejati pada $K$ (sebab $d(x, F) = 0$ berarti $x \in \bar F =
F$); dan pada $K$ yang \omterm{def:b3:topology:compact}{kompak} ia mencapai minimum $m > 0$: jadi $d(K, F)
\geq m$. Contoh penyangkal tanpa kekompakan: $F_1 = \{n : n \geq
2\}$ dan $F_2 = \{n + \frac1n : n \geq 2\}$ tertutup,
saling lepas, tetapi berjarak $\inf_n \frac1n = 0$.

(c) Jika $X$ \omterm{def:b3:topology:compact}{kompak}, setiap $f$ yang \omterm{def:b3:topology:continuity}{kontinu} bersifat terbatas
(\cref{thm:b3:topology:compactprops}(2)). Sebaliknya, jika $X$
tidak \omterm{def:b3:topology:compact}{kompak}, ambillah $(x_n)$ yang tak punya subbarisan konvergen
(\cref{thm:b3:topology:metriccompact}); dengan beralih ke
subbarisan kita boleh menganggap $x_n$ berbeda berpasangan, dan tak ada
titik $X$ yang menjadi titik kumpul barisan itu, sehingga
\[
r_n = \min\Bigl(2^{-n},\ \tfrac13\,d\bigl(x_n, \{x_m : m \neq
n\}\bigr)\Bigr) > 0,
\]
dan bola $B(x_n, r_n)$ saling lepas berpasangan (sebab titik persekutuan
bola ke-$n$ dan ke-$m$ akan memberikan $d(x_n, x_m) < r_n +
r_m \leq \frac23 d(x_n, x_m)$). Definisikan
\[
f(x) = \sum_{n} n\,\max\Bigl(0,\ 1 - \frac{d(x,
x_n)}{r_n}\Bigr):
\]
Pada setiap $x$ paling banyak satu sukunya tak nol, dan $f(x_n) = n$.
\omterm{def:b3:topology:continuity}{Kekontinuan} di $x$: misalkan $y_k \to x$; setiap nilai $f(y_k)$
entah $0$ entah sebuah nilai gundukan tunggal $\mathrm{bump}_{n_k}(y_k)$.
Jika suatu indeks $n$ muncul tak berhingga sering, maka sepanjang
subbarisan itu $f(y_k) = \mathrm{bump}_n(y_k) \to
\mathrm{bump}_n(x)$, yang sama dengan $f(x)$ (sebab jika $x \in B(x_n,
r_n)$, seluruh ekornya terletak di bola terbuka itu dan tak ada indeks
lain yang muncul; sedangkan jika $x \notin B(x_n, r_n)$ maka
$\mathrm{bump}_n(x) = 0 = f(x)$, karena $x$, yang melekat pada bola
ke-$n$, tak terletak di bola terbuka lain mana pun). Jika $n_k \to \infty$:
maka $d(y_k, x_{n_k}) < r_{n_k} \leq 2^{-n_k} \to 0$, sehingga $x_{n_k}
\to x$, dan itu menjadikan $x$ titik kumpul $(x_n)$ --- yang sudah dikecualikan; jadi
kasus ini hanya menyangkut berhingga banyak $k$. Pada setiap kasus
$f(y_k) \to f(x)$: jadi $f$ \omterm{def:b3:topology:continuity}{kontinu}, sekaligus tak terbatas.
\end{solution}

\begin{solution}{exo:b3:topology:7}
Andaikan tak ada $\delta$ yang berhasil: untuk setiap $n$ ada $A_n$ berdiameter
$< 1/n$ yang tak termuat di satu $U_i$ pun; pilihlah $x_n \in
A_n$. Menurut kekompakannya (\cref{thm:b3:topology:metriccompact}),
ada subbarisan $x_{n_k} \to x$; pilihlah $i$ dan $r > 0$ dengan $B(x,
r) \subseteq U_i$. Untuk $k$ besar: $d(x_{n_k}, x) < r/2$ dan
$1/n_k < r/2$, sehingga $A_{n_k} \subseteq B(x_{n_k}, 1/n_k)
\subseteq B(x, r) \subseteq U_i$ --- kontradiksi. Untuk Heine:
diberikan $\varepsilon$, selimuti $Y$ dengan bola $B(y, \varepsilon/2)$;
prapetanya membentuk selimut terbuka bagi $X$; misalkan $\delta$ sebuah
bilangan Lebesgue: jika $d(x, x') < \delta$, maka pasangan $\{x, x'\}$
berdiameter $< \delta$, terletak pada satu prapeta, sehingga $d(f(x),
f(x')) < \varepsilon$.
\end{solution}

\begin{solution}{exo:b3:topology:8}
(a) $GL_n(\R) = \det^{-1}(\R\setminus\{0\})$ terbuka (sebab $\det$
polinomial, jadi \omterm{def:b3:topology:continuity}{kontinu}). Ketakterhubungannya: $\det$ memetakannya pada
$\R\setminus\{0\}$, sedangkan \omterm{def:b3:topology:connected}{ruang terhubung} berpeta \omterm{def:b3:topology:continuity}{kontinu} yang
\omterm{def:b3:topology:connected}{terhubung} (\cref{thm:b3:topology:connectedbasics}(2)); padahal
$\R\setminus\{0\}$ bukan selang. Untuk $GL_n(\C)$: untuk $A, B$
yang terbalikkan, $p(\lambda) = \det\bigl((1-\lambda)A + \lambda
B\bigr)$ merupakan polinomial dalam $\lambda$ dengan $p(0), p(1) \neq 0$,
sehingga tak nol secara identik: jadi ia berakar berhingga banyak di $\C$.
Bidang dikurangi berhingga banyak titik bersifat \omterm{def:b3:topology:pathconnected}{terhubung lintasan} (kelilingi
titik itu), sehingga ada \omterm{def:b3:topology:pathconnected}{lintasan} $\lambda(t)$
dari $0$ ke $1$ dengan $p(\lambda(t)) \neq 0$: jadi $t \mapsto
(1-\lambda(t))A + \lambda(t)B$ merupakan \omterm{def:b3:topology:pathconnected}{lintasan} di $GL_n(\C)$.

(b) Determinan $\det(A + \varepsilon I) = \chi_A(-\varepsilon)\cdot(\pm1)$
lenyap untuk paling banyak $n$ nilai $\varepsilon$: sehingga matriks terbalikkan
$A + \varepsilon I \to A$ ketika $\varepsilon \to 0$: itulah
kepadatannya.
\end{solution}

\begin{solution}{exo:b3:topology:9}
(a) Misalkan $A \subseteq \Q$ memuat $p < q$ lalu pilihlah bilangan irasional
$\alpha \in (p, q)$: maka $A = (A \cap (-\infty, \alpha)) \sqcup
(A\cap(\alpha, +\infty))$ memecah $A$ menjadi dua kepingan tak kosong yang
terbuka relatif: jadi $A$ tak \omterm{def:b3:topology:connected}{terhubung}. Komponennya beranggota tunggal.
Kekompakan lokalnya gagal: \omterm{def:b3:topology:topology}{persekitaran} \omterm{def:b3:topology:compact}{kompak} $V$ dari $0$ di
$\Q$ memuat $[-r, r]\cap\Q$ untuk suatu $r > 0$, yang tertutup
di $V$, sehingga \omterm{def:b3:topology:compact}{kompak}; padahal barisan bilangan rasional di $[-r,r]$ yang
konvergen (di $\R$) ke bilangan irasional tak punya subbarisan yang
konvergen \emph{di $\Q$}: jadi bertentangan dengan
\cref{thm:b3:topology:metriccompact}.

(b) Misalkan $\gamma = (\gamma_1, \gamma_2) \colon [0,1] \to S$ sebuah
\omterm{def:b3:topology:pathconnected}{lintasan} dengan $\gamma(0) = (0,0)$ dan $\gamma(1) \in \Gamma$.
Himpunan $\{t : \gamma_1(t) = 0\}$ tertutup dan tak memuat
$1$; misalkan $t^*$ supremumnya, sehingga $\gamma_1(t^*) = 0$ dan
$\gamma_1 > 0$ pada $(t^*, 1]$. Menurut \omterm{def:b3:topology:continuity}{kekontinuan} $\gamma_2$ di
$t^*$, pilihlah $\delta$ dengan $\abs{\gamma_2(t) -
\gamma_2(t^*)} < \tfrac12$ untuk $t \in [t^*, t^* + \delta]$. Tetapkan
$t_0 = t^* + \delta$ (boleh dianggap $\leq 1$): maka $\gamma_1(t_0) > 0$,
dan $\gamma_1$ mengambil setiap nilai $(0, \gamma_1(t_0)]$ pada
$[t^*, t_0]$ (menurut teorema nilai antara,
\cref{thm:b3:topology:connectedbasics}). Pilihlah $k$ besar dengan
$a_k = \frac1{\pi/2 + 2k\pi} < \gamma_1(t_0)$ dan $b_k =
\frac1{3\pi/2 + 2k\pi} < \gamma_1(t_0)$: maka ada $s, s' \in
(t^*, t_0]$ dengan $\gamma_1(s) = a_k$, $\gamma_1(s') = b_k$;
dan karena titiknya terletak pada grafiknya, $\gamma_2(s) = \sin(1/a_k) =
1$ serta $\gamma_2(s') = -1$. Keduanya tak mungkin berjarak $\frac12$ dari
$\gamma_2(t^*)$: kontradiksi. Jadi $S$ \omterm{def:b3:topology:connected}{terhubung} tetapi tidak
\omterm{def:b3:topology:pathconnected}{terhubung lintasan}.
\end{solution}

\begin{solution}{exo:b3:topology:10}
(a) Himpunan $C_n$ tepat terdiri atas $x$ yang mempunyai ekspansi
terner dengan angka di $\{0, 2\}$ sampai peringkat $n$ (lewat induksi:
membuang sepertiga tengah membuang angka pertama $1$, dan seterusnya; sedangkan ujungnya
mempunyai dua ekspansi, salah satunya tanpa angka $1$), sehingga $C = \bigcap C_n$
merupakan himpunan jumlah $\sum a_n3^{-n}$ dengan $a_n \in \{0,2\}$. Kekompakannya:
setiap $C_n$ berupa gabungan berhingga selang tertutup; jadi $C$ tertutup
di $[0,1]$. \omterm{def:b3:topology:interior}{Interiornya} kosong: sebab $C \subseteq C_n$, yaitu gabungan
selang berpanjang $3^{-n}$; dan selang \omterm{def:b3:topology:interior}{interior} berpanjang
$\ell > 0$ harus muat di salah satunya untuk setiap $n$. Kesempurnaannya: diberikan
$x \in C$ dan $m$, baliklah angka $a_m$: titik barunya berada di
$C$, berbeda, dan berjarak dalam $2\cdot3^{-m}$ dari $x$.

(b) Jika $x \neq y \in C$, maka ekspansinya pertama kali berbeda pada suatu
peringkat $k$; di antara keduanya terletak sebuah selang sepertiga tengah yang dibuang
(yakni celah pada peringkat $k$ yang memisahkan angka $0$ dari angka $2$),
sehingga memberikan titik $c \notin C$ dengan $x < c < y$ (katakanlah): jadi $C =
(C \cap (-\infty, c)) \sqcup (C \cap (c, \infty))$ memecah sembarang
himpunan bagian yang memuat kedua titik itu. Jadi komponennya beranggota tunggal.

(c) Pemetaan angka $\beta\colon x \mapsto (a_n(x))$ merupakan
bijeksi pada $\{0,2\}^\N$ (lewat keberadaan dan ketunggalan
ekspansi-$\{0,2\}$: barisan angka yang berbeda memberikan titik yang
berjarak $\geq 3^{-k} > 0$ pada tempat pertama perbedaannya, seperti pada
\cref{pb:b3:topology:1}, pertanyaan 1). Ia \omterm{def:b3:topology:continuity}{kontinu}: sebab
$\abs{x - y} < 3^{-m}$ memaksa $m$ angka pertamanya bersesuaian
(dengan perhitungan yang sama), sehingga $\beta$ memetakan bola kecil ke dalam kotak
basisnya. Dan bijeksi \omterm{def:b3:topology:continuity}{kontinu} dari $C$ yang \omterm{def:b3:topology:compact}{kompak} ke hasil kali yang
\omterm{def:b3:topology:hausdorff}{Hausdorff} merupakan \omterm{def:b3:topology:continuity}{homeomorfisma}
(\cref{cor:b3:topology:compacthomeo}; hasil kalinya \omterm{def:b3:topology:hausdorff}{Hausdorff}:
pisahkan pada koordinat yang berbeda). Ketakterbilangannya: lewat diagonal
Cantor pada $\{0,2\}^\N$. Akhirnya $\{0,2\}^\N \times
\{0,2\}^\N \cong \{0,2\}^\N$ lewat penyelipan angka (yang merupakan
\omterm{def:b3:topology:continuity}{homeomorfisma}: \omterm{def:b3:topology:continuity}{kontinu} komponen demi komponen pada kedua arahnya), sehingga $C
\times C \cong C$.
\end{solution}

\begin{solution}{exo:b3:topology:11}
Untuk $y \in \R^n$, tetapkan
\[
\tau(y) = \Bigl(\frac{2y}{\norm y^2 + 1},\ \frac{\norm y^2 -
1}{\norm y^2 + 1}\Bigr) \in \R^n\times\R = \R^{n+1}.
\]
Mudah diperiksa $\norm{\tau(y)} = 1$, $\tau(y) \neq N$, dan
$\sigma(\tau(y)) = y$, $\tau(\sigma(x)) = x$: jadi $\sigma$ dan
$\tau$ saling invers, keduanya \omterm{def:b3:topology:continuity}{kontinu} (karena rumus rasional
yang penyebutnya tak lenyap): sehingga $S^n\setminus\{N\} \cong
\R^n$. Perluaslah menjadi $\hat\sigma \colon S^n \to \widehat{\R^n}$ dengan
$N \mapsto \infty$: sebuah bijeksi yang \omterm{def:b3:topology:continuity}{kontinu} di setiap titik
$S^n\setminus\{N\}$, dan juga di $N$: sebab \omterm{def:b3:topology:topology}{persekitaran} \omterm{def:b3:topology:basis}{basis}
$\infty$ berupa $\widehat{\R^n}\setminus K$ dengan $K$ \omterm{def:b3:topology:compact}{kompak} dan $K
\subseteq \bar B(0, R)$; lalu karena $\norm{\sigma(x)}^2 = \frac{1 +
x_{n+1}}{1 - x_{n+1}} \to \infty$ ketika $x_{n+1} \to 1$, himpunan
$\{x \in S^n : x_{n+1} > 1 - \eta\}$ terpetakan di luar $\bar B(0,R)$
untuk $\eta$ kecil: itulah \omterm{def:b3:topology:continuity}{kekontinuannya}. Dan bijeksi \omterm{def:b3:topology:continuity}{kontinu} dari
$S^n$ yang \omterm{def:b3:topology:compact}{kompak} ke $\widehat{\R^n}$ yang \omterm{def:b3:topology:hausdorff}{Hausdorff} merupakan
\omterm{def:b3:topology:continuity}{homeomorfisma}. Untuk membuang titik lain $p$: sebuah rotasi $S^n$
memetakan $p$ ke $N$ (rotasi merupakan \omterm{def:b3:topology:continuity}{homeomorfisma}), sehingga kembali ke
kasus yang sudah dihitung: $S^n \setminus \{p\} \cong \R^n$.
\end{solution}

\begin{solution}{exo:b3:topology:12}
(a) Himpunan $T$ terbatas, dan tertutup: sebab limit titik $T$
yang absisnya $\to x_0 > 0$ tetap berada pada grafiknya (yang tertutup
lokal) menurut \omterm{def:b3:topology:continuity}{kekontinuan} $\sin\frac1x$ pada $\intoc01$; sedangkan limit
yang absisnya $\to 0$ berordinat di $\intcc{-1}1$, sehingga
terletak pada ruas garisnya. Ia \omterm{def:b3:topology:compact}{kompak} menurut Heine--Borel
(\cref{cor:b3:topology:heineborel}). Sebagai \omterm{def:b3:topology:interior}{penutup} grafik
$\Gamma$: setiap titik $(0, y)$ dengan $y \in \intcc{-1}1$ merupakan
limit titik grafiknya --- selesaikan $\sin\frac1x = y$ di dekat $0$
(sebab fungsinya menyapu $\intcc{-1}1$ pada setiap selang
$[\frac1{(k+1)\pi}, \frac1{k\pi}]$): jadi $\bar\Gamma = T$.

(b) Himpunan $\Gamma$ merupakan peta \omterm{def:b3:topology:continuity}{kontinu} $\intoc01$ yang
\omterm{def:b3:topology:connected}{terhubung} di bawah $x \mapsto (x, \sin\frac1x)$: jadi \omterm{def:b3:topology:connected}{terhubung};
dan \omterm{def:b3:topology:interior}{penutup} himpunan \omterm{def:b3:topology:connected}{terhubung} tetap \omterm{def:b3:topology:connected}{terhubung}
(\cref{thm:b3:topology:connectedbasics}): sehingga $T = \bar\Gamma$
\omterm{def:b3:topology:connected}{terhubung}.

(c) Andaikan $\gamma\colon\intcc01\to T$ \omterm{def:b3:topology:continuity}{kontinu} dengan
$\gamma(0) = (0,0)$, $\gamma(1) = (1, \sin1)$, lalu misalkan $t^*
= \sup\{t : \gamma_1(t) = 0\}$: menurut \omterm{def:b3:topology:continuity}{kekontinuannya} $\gamma_1(t^*)
= 0$ dan $\gamma_1 > 0$ pada $\intoc{t^*}1$. Untuk setiap
$\delta > 0$, pada selang $\intoc{t^*}{t^*+\delta}$ nilai
$\gamma_1$ mengambil semua nilai di suatu $\intoc0\eta$
(menurut teorema nilai antara, sebab $\gamma_1(t^* + \delta) > 0$);
khususnya ia memuat absis berbentuk
$\frac1{\pi/2 + 2k\pi}$ dan $\frac1{-\pi/2 + 2k\pi}$ untuk
$k$ sebesar apa pun, dan di situ $\gamma_2 = \pm1$. Jadi pada
setiap \omterm{def:b3:topology:topology}{persekitaran} kanan $t^*$, $\gamma_2$ mengambil kedua
nilai $1$ dan $-1$: sehingga $\gamma_2$ tak berlimit di $t^{*+}$,
dan itu bertentangan dengan \omterm{def:b3:topology:continuity}{kekontinuannya}. Jadi tak ada \omterm{def:b3:topology:pathconnected}{lintasan} yang demikian.

(d) Himpunan $T$ \omterm{def:b3:topology:connected}{terhubung} tetapi tidak \omterm{def:b3:topology:pathconnected}{terhubung lintasan}: jadi kedua gagasan itu
berbeda. Untuk $U \subseteq \R^n$ yang terbuka dan \omterm{def:b3:topology:connected}{terhubung} serta $a \in U$,
misalkan $V$ himpunan titik $U$ yang dapat dihubungkan ke $a$ oleh
sebuah \omterm{def:b3:topology:pathconnected}{lintasan} di $U$. Himpunan $V$ terbuka: sebab di sekitar $x \in V$ ada bola $B(x, r)
\subseteq U$ yang berbentuk bintang, dan menyambung \omterm{def:b3:topology:pathconnected}{lintasan} menuju $x$
dengan sebuah ruas garis mencapai setiap titik bola itu. Himpunan $V$ juga
tertutup di $U$: sebab jika $x \in U \setminus V$, argumen bola yang sama
menunjukkan $B(x, r) \cap V = \varnothing$ (sebab titik
$V$ di dalam bola itu akan menghubungkan $a$ ke $x$). Karena tak kosong ($a \in
V$), terbuka dan tertutup di $U$ yang \omterm{def:b3:topology:connected}{terhubung}: maka $V = U$. Kurva
sinus itu mengelak dari hal ini dengan menjadi tertutup dan berinterior kosong: sebab
titik ``buruknya'' $(0,0)$ tak punya bola di dalam $T$ untuk menjembatani
ayunannya.
\end{solution}

\begin{solution}{pb:b3:topology:1}
\textbf{1.} Andaikan ekspansi $x, y \in C$ pertama kali berbeda
pada peringkat $k \leq m$, katakanlah $b_k(x) = 1$, $b_k(y) = 0$. Maka
\[
x - y \;\geq\; \frac{2}{3^k} - \sum_{j > k}\frac{2}{3^j}
= \frac{2}{3^k} - \frac{1}{3^k} = \frac1{3^k} \geq \frac1{3^m}:
\]
kontraposisinya: $\abs{x - y} < 3^{-m}$ memaksa kesesuaian sampai
peringkat $m$.

\textbf{2.} Menurut pertanyaan 1, $b_n$ konstan pada $C \cap (x -
3^{-n}, x + 3^{-n})$: jadi konstan lokal, sehingga \omterm{def:b3:topology:continuity}{kontinu}.
Pemetaan $x \mapsto (b_n(x))_n \in \{0,1\}^\N$ bersifat \omterm{def:b3:topology:continuity}{kontinu}
(komponen demi komponen, \cref{prop:b3:topology:universal}(a)) dan
bijektif (karena ekspansi angka-$\{0,2\}$ tunggal); dan dari $C$ yang
\omterm{def:b3:topology:compact}{kompak} ke \omterm{def:b3:topology:hausdorff}{ruang Hausdorff}, ia merupakan \omterm{def:b3:topology:continuity}{homeomorfisma}
(\cref{cor:b3:topology:compacthomeo}).

\textbf{3.} \omterm{def:b3:topology:continuity}{Kekontinuannya}: jika $\abs{x - y} < 3^{-m}$, maka $m$ angka
pertamanya bersesuaian, sehingga $\abs{g(x) - g(y)} \leq \sum_{n > m}2^{-n}
= 2^{-m}$. Kesurjektifannya: setiap $t \in [0,1]$ mempunyai ekspansi
biner $t = \sum b_n2^{-n}$, dan $t = g\bigl(\sum
2b_n3^{-n}\bigr)$. Tidak injektif: sebab $g$ mengenali kedua titik
Cantor berangka $(0,1,1,1,\dots)$ dan $(1,0,0,\dots)$ ---
keduanya terpetakan ke $\frac12$ (lewat kedwiartian diadik $0.0111\ldots_2 =
0.1000\ldots_2$).

\textbf{4.} Setiap komponen $\Phi$ \omterm{def:b3:topology:continuity}{kontinu} lewat taksiran yang sama
(sebab angkanya merupakan subbarisan $b_n$).
Kesurjektifannya: diberikan $(s, t) \in [0,1]^2$, pilihlah angka biner
$(c_n)$ dari $s$ dan $(d_n)$ dari $t$, lalu selipkan berselang: titik $x
\in C$ dengan $b_{2n-1}(x) = c_n$, $b_{2n}(x) = d_n$ memenuhi
$\Phi(x) = (s,t)$.

\textbf{5.} Celahnya saling lepas berpasangan dengan ujung di
$C$, sehingga rumusnya mendefinisikan $f$ tanpa kedwiartian pada $[0,1]$; dan di
ujung sebuah celah rumus afinnya mengembalikan $\Phi(u)$ dan
$\Phi(v)$: jadi konsisten dengan $f = \Phi$ pada $C$. Kesurjektifannya:
sudah $f(C) = \Phi(C) = [0,1]^2$.

\textbf{6.} Di $t_0 \notin C$: titik $t_0$ terletak pada celah terbuka tempat
$f$ bersifat afin: jadi \omterm{def:b3:topology:continuity}{kontinu}. Di $x \in C$: kita catat lebih dahulu
dua taksiran.

\emph{(i) Pada $C$}: jika $\abs{y - x} \leq 3^{-M}$ dengan $y \in C$,
maka angkanya bersesuaian sampai $M$, sehingga setiap komponen $\Phi(y)
- \Phi(x)$ paling besar $\sum_{n > \lfloor M/2\rfloor} 2^{-n} =
2^{-\lfloor M/2\rfloor}$.

\emph{(ii) Menyeberangi sebuah celah}: celah $(u,v)$ yang dibuang pada tahap $k$
berpanjang $3^{-k}$, dan ujungnya berangka sesuai sampai
peringkat $k - 1$ (sebab keduanya berbeda mulai peringkat $k$), sehingga
$\norm{\Phi(v) - \Phi(u)}_\infty \leq
2^{-\lfloor (k-1)/2\rfloor}$.

Sekarang misalkan $\abs{t - x} < 3^{-M}$ dengan $t$ pada sebuah celah $(u, v)$ dari
tahap $k$, dan katakanlah $u$ ujung di pihak $x$, sehingga
$\abs{u - x} < 3^{-M}$ dan $0 < t - u < 3^{-M}$. Maka
\[
\norm{f(t) - \Phi(x)}_\infty
\leq \norm{\Phi(u) - \Phi(x)}_\infty
+ \frac{t - u}{v - u}\,\norm{\Phi(v) - \Phi(u)}_\infty
\leq 2^{-\lfloor M/2\rfloor} + \min\bigl(1,\ 3^{k - M}\bigr)\,
2^{-\lfloor(k-1)/2\rfloor} .
\]
Untuk $k \geq M$ suku terakhirnya $\leq 2^{-\lfloor(M-1)/2
\rfloor}$; sedangkan untuk $k < M$, karena $2^{-\lfloor(k-1)/2\rfloor} \leq
\sqrt2\,\cdot 2^{-k/2}$,
\[
3^{k-M}\,2^{-\lfloor(k-1)/2\rfloor}
\leq \sqrt2\,\Bigl(\frac{3}{\sqrt2}\Bigr)^{k} 3^{-M}
\leq \sqrt2\,\Bigl(\frac{3}{\sqrt2}\Bigr)^{M} 3^{-M}
= \sqrt2\; 2^{-M/2}
\]
(sebab kuantitas di tengahnya naik terhadap $k$). Pada semua kasus
$\norm{f(t) - \Phi(x)}_\infty \leq C\,2^{-M/2}$
dengan $C$ berupa konstanta mutlak: jadi membiarkan $M \to \infty$ membuktikan
\omterm{def:b3:topology:continuity}{kekontinuannya} di $x$ (sedangkan titik $t \in C$ tercakup oleh (i)). Karena itu
$f$ merupakan surjeksi \omterm{def:b3:topology:continuity}{kontinu} $[0,1] \to [0,1]^2$ --- bahkan
taksirannya menunjukkan ia Hölder bercitarasa eksponen $\log 2/(2\log
3)$, tetapi kekontinuanlah satu-satunya yang kita klaim.

\textbf{7.} Untuk $[0,1]^n$: pecahlah angka $x \in C$ menjadi
$n$ subbarisan berselang lalu ulangi pertanyaan 4--6
kata demi kata. Untuk $\R \to \R^2$: misalkan $h \colon [0,1] \to [-1,1]^2$
sebuah surjeksi \omterm{def:b3:topology:continuity}{kontinu} (skalakan kembali $f$). Definisikan $F$ pada $[m, m
+ \frac12]$ (dengan $m \in \Z$) sebagai salinan $h$ yang diskalakan untuk mengisi
$[-(m+1), m+1]^2$ bila $m \geq 0$ (dan secara simetris bila $m <
0$), lalu pada $[m + \frac12, m + 1]$ sebagai ruas afin yang
menghubungkan nilai di ujungnya: maka $F$ \omterm{def:b3:topology:continuity}{kontinu} (lewat perekatan pada
kepingan tertutup, \cref{prop:b3:topology:gluing}(b)) dan petanya
memuat $\bigcup_m [-(m+1), m+1]^2 = \R^2$.

\textbf{8.} Ruang $[0,1]$ \omterm{def:b3:topology:compact}{kompak} dan $[0,1]^2$ \omterm{def:b3:topology:hausdorff}{Hausdorff}: sehingga
injeksi \omterm{def:b3:topology:continuity}{kontinu} merupakan \omterm{def:b3:topology:continuity}{homeomorfisma} pada petanya
(\cref{cor:b3:topology:compacthomeo} yang diterapkan pada korestriksinya).

\textbf{9.} Buanglah $\frac12$: maka $[0,1]\setminus\{\frac12\}$ menjadi
tak \omterm{def:b3:topology:connected}{terhubung}. Seandainya $h \colon [0,1] \to [0,1]^2$ sebuah
\omterm{def:b3:topology:continuity}{homeomorfisma}, maka $[0,1]^2\setminus\{h(\frac12)\}$ juga akan
tak \omterm{def:b3:topology:connected}{terhubung} (karena peta \omterm{def:b3:topology:continuity}{homeomorfik} ruang tak \omterm{def:b3:topology:connected}{terhubung}
tetap tak \omterm{def:b3:topology:connected}{terhubung}); padahal persegi dikurangi sebuah titik bersifat
\omterm{def:b3:topology:pathconnected}{terhubung lintasan}: hubungkan dua titiknya dengan \omterm{def:b3:topology:pathconnected}{lintasan} dua ruas yang menghindari
lubangnya. Kontradiksi: jadi $[0,1] \not\cong [0,1]^2$.

\textbf{10.} Menurut pertanyaan 8, surjeksi \omterm{def:b3:topology:continuity}{kontinu} yang injektif
$[0,1] \to [0,1]^2$ akan menjadi \omterm{def:b3:topology:continuity}{homeomorfisma}, dan itu bertentangan dengan
pertanyaan 9. Kekompakannya masuk persis pada
\cref{cor:b3:topology:compacthomeo}: sebab di situlah invers sebuah
bijeksi \omterm{def:b3:topology:continuity}{kontinu} menjadi \omterm{def:b3:topology:continuity}{kontinu} (karena peta himpunan tertutup
bersifat \omterm{def:b3:topology:compact}{kompak}, sehingga tertutup). Tanpa kekompakan
kesimpulannya sungguh gagal: $[0, 2\pi) \to S^1$ merupakan bijeksi
\omterm{def:b3:topology:continuity}{kontinu} yang bukan \omterm{def:b3:topology:continuity}{homeomorfisma}.

\textbf{11.} Yang dibuktikan di sini: \emph{ada surjeksi
\omterm{def:b3:topology:continuity}{kontinu}, tetapi tak ada bijeksi \omterm{def:b3:topology:continuity}{kontinu}, dari $[0,1]$ pada
$[0,1]^2$; khususnya $[0,1] \not\cong [0,1]^2$.} Jadi
``dimensi'' tidak terpelihara oleh surjeksi \omterm{def:b3:topology:continuity}{kontinu} ---
kardinalitas dan bahkan \omterm{def:b3:topology:continuity}{kekontinuan} tak dapat melihatnya --- tetapi ia
\emph{memang} invarian \omterm{def:b3:topology:topology}{topologi} pada tingkat
\omterm{def:b3:topology:continuity}{homeomorfisma}, setidaknya untuk dimensi $1$ lawan $2$, tempat
keterhubungan setelah penghapusan sudah mencukupi. Pernyataan umumnya
--- yaitu \emph{invariansi domain}: bahwa $\R^m \cong \R^n$ mengakibatkan $m =
n$, dan bahwa injeksi \omterm{def:b3:topology:continuity}{kontinu} $\R^n \to \R^n$ bersifat terbuka ---
memang benar tetapi menuntut \omterm{def:b3:topology:topology}{topologi} aljabar (homologi), yang melampaui
kuliah ini.

\textbf{12.} Dalam terner, $\frac C2 = \{\sum_na_n3^{-n} : a_n
\in \{0, 1\}\}$ (karena membagi dua angka $\{0,2\}$ memberikan angka
$\{0,1\}$). Untuk $t \in [0,1]$, ambillah sembarang ekspansi terner $t =
\sum_nd_n3^{-n}$ dengan $d_n \in \{0,1,2\}$, lalu pecahlah setiap angkanya
sebagai $d_n = a_n + b_n$ dengan $a_n, b_n \in \{0,1\}$ ($0 = 0{+}0$,
$1 = 1{+}0$, $2 = 1{+}1$): maka $t \in \frac C2 + \frac C2$.
Intinya adalah bahwa setiap angkanya terpecah \emph{di dalam} $\{0,1\}$, sehingga
tak ada simpanan yang merambat dan angkanya dapat digarap
secara bebas. Inklusi sebaliknya jelas (sebab jumlah angkanya tetap
$\leq 2$). Setelah diskalakan dengan $2$: $C + C = [0,2]$.

\textbf{13.} Pemetaan $(x, y) \mapsto x + y$ mengirim $C \times C$
pada $C + C = [0,2]$: jadi debu itu, yang tak memuat persegi mana pun
(karena \omterm{def:b3:topology:interior}{interiornya} kosong: \cref{exo:b3:topology:10}), terproyeksikan sepanjang
antidiagonalnya pada ruas garis penuh. Keserupaan dirinya: sebab
angka terner pertama sebuah titik Cantor bernilai $0$ atau $2$, dan
melucutinya memberikan $C = \frac C3 \cup \bigl(\frac C3 +
\frac23\bigr)$ --- yaitu persamaan titik tetap yang membangkitkan
setiap gambar $C$.

\textbf{14.} Melalui $C \cong \{0,1\}^\N$ (pertanyaan 2),
penyelipan $(\varepsilon, \varepsilon') \mapsto
(\varepsilon_1, \varepsilon_1', \varepsilon_2,
\varepsilon_2', \dots)$ merupakan bijeksi
$\{0,1\}^\N\times\{0,1\}^\N \to \{0,1\}^\N$ yang \omterm{def:b3:topology:continuity}{kontinu} pada kedua
arahnya (sebab setiap koordinat keluarannya bergantung pada satu koordinat
masukan saja; lihat \omterm{def:b3:topology:constructions}{topologi hasil kalinya}). Karena itu $C\times C \cong C$ dan
secara induktif $C^n \cong C$. Ruang yang akrab tak berbagi sifat ini:
$\R^2 \not\cong \R$ dan $[0,1]^2 \not\cong [0,1]$
(pertanyaan 11); sedangkan hasil kali tak berhingga memang berbagi: $([0,1]^\N)^2 \cong
[0,1]^\N$ lewat penyelipan yang sama.

\textbf{15.} $\sum_{k\geq1}2\cdot3^{-2k} = 2\cdot\frac{1/9}{1
- 1/9} = \frac14$: jadi ekspansi terner $\frac14$ adalah
$0.\overline{02}$, dengan angka di $\{0,2\}$, sehingga $\frac14 \in
C$; dan karena ia bukan ekspansi berhingga maupun ekor yang
akhirnya bernilai $2$, ia bukan ujung sepertiga tengah mana pun yang dibuang. Ujungnya
membentuk himpunan terbilang (dua per selang yang dibuang, dan selangnya
terbilang banyaknya), sedangkan $C \cong \{0,1\}^\N$ tak terbilang
(lewat argumen diagonal, atau pertanyaan 19): jadi hampir setiap titik
Cantor, seperti $\frac14$, tak terlihat pada gambar ujung yang
biasa.

\textbf{16.} Kesamaan $g(x) = g(y)$ berarti barisan biner
$(b_n(x))$ dan $(b_n(y))$ mewakili bilangan real yang sama. Dan barisan
biner berbeda yang mewakili bilangan sama tepat muncul pada
kedwiartian diadik $0.w01^\infty = 0.w10^\infty$: jadi paling banyak
dua prapeta, dan tepat dua persis pada bilangan
rasional diadik di $(0,1)$. Jadi $g$ meruntuhkan $C$ pada $[0,1]$ dengan
merekatkan terbilang banyak pasangan --- yaitu kerangka
kombinatorik tangga Cantor pada \cref{ch:b3:measure}.

\textbf{17.} Titik $x$ tak terisolasi di $P$, sehingga $B(x, r/2)$
memuat suatu $y_0 \in P \setminus \{x\}$; dan $y_0$ pun tak
terisolasi, sehingga $B(x, r/2)$ memuat titik kedua $y_1
\in P$. Sembarang $\rho \leq \min\bigl(\frac{d(y_0,y_1)}3,
\frac r4\bigr)$ membuat $\bar B(y_0, \rho)$ dan $\bar B(y_1,
\rho)$ saling lepas dan termuat di $B(x, r)$. Keduanya berpusat
di titik $P$, dan di situ penarikan dua titik yang sama dapat
diulang: jadi kesempurnaan merupakan pasokan tak habis-habis titik terdekat
yang menjaga rekursinya tetap hidup selamanya.

\textbf{18.} Bangunlah $F_w$ dengan induksi pada panjang katanya:
$F_\varnothing = P$, dan di dalam setiap bola $B(y_w, \rho_w)$
pertanyaan 17 melahirkan dua bola tertutup saling lepas berjari-jari $\leq
\min(\rho_w/2,\ 2^{-\abs w-1}\operatorname{diam}P)$ yang berpusat
di titik $P$; lalu tetapkan $F_{wi} = \bar B(y_{wi}, \rho_{wi})
\cap P$: tak kosong (memuat pusatnya) dan \omterm{def:b3:topology:compact}{kompak}. Untuk
$\varepsilon$ yang tak berhingga, \omterm{def:b3:topology:compact}{kompak} tak kosong yang bersarang
$F_{\varepsilon_1\cdots\varepsilon_m}$ beririsan tak kosong (lewat sifat irisan
berhingga pada $P$ yang \omterm{def:b3:topology:compact}{kompak}), dengan diameter $0$: jadi satu titik $h(\varepsilon)$.

\textbf{19.} Kata yang pertama kali berbeda pada peringkat $m$ mengirim
petanya ke dua himpunan yang \emph{saling lepas} $F_{w0}, F_{w1}$
(dengan awalan bersama $w$): jadi $h$ injektif. Jika dua kata bersesuaian sampai
peringkat $m$, maka kedua petanya terletak di himpunan berdiameter $\leq
2^{-m}\operatorname{diam}P$: jadi $h$ \omterm{def:b3:topology:continuity}{kontinu}. Karena itu
$\{0,1\}^\N$ yang tak terbilang terinjeksi ke $P$: jadi setiap ruang metrik \omterm{def:b3:topology:compact}{kompak}
yang \omterm{prop:b3:galois:perfect}{sempurna} bersifat tak terbilang --- termasuk $[0,1]$ dan $C$.

\textbf{20.} Pemetaan $h$ merupakan injeksi \omterm{def:b3:topology:continuity}{kontinu} dari $\{0,1\}^\N$
yang \omterm{def:b3:topology:compact}{kompak} ke $P$ yang \omterm{def:b3:topology:hausdorff}{Hausdorff} (karena metrik):
\cref{cor:b3:topology:compacthomeo} meningkatkannya menjadi
\omterm{def:b3:topology:continuity}{homeomorfisma} pada petanya; dan $\{0,1\}^\N \cong C$
(pertanyaan 2). Jadi setiap ruang metrik \omterm{def:b3:topology:compact}{kompak} yang \omterm{prop:b3:galois:perfect}{sempurna} memuat
salinan himpunan Cantor: kesempurnaan mempunyai benih universal, dan
benih itu milik Cantor.

\textbf{21.} Ruang metrik \omterm{def:b3:topology:compact}{kompak} yang terbilang tak mungkin
\omterm{prop:b3:galois:perfect}{sempurna} (pertanyaan 19), sehingga ia mempunyai titik terisolasi. Di $K =
\{0\}\cup\{\frac1n : n \geq 1\}$ setiap titik $\frac1n$ bersifat
terisolasi sedangkan $0$ tidak: jadi titik terisolasi bahkan dapat \omterm{def:b3:topology:interior}{padat}
tanpa membuat ruangnya diskret --- sebab kekompakan menahan
limitnya tetap di dalam.

\textbf{22.} Misalkan $P$ himpunan titik kondensasi $F$
dan $\mathcal B$ keluarga terbilang berisi selang terbuka
berujung rasional. Setiap titik $F$ yang bukan kondensasi terletak
di suatu $I \in \mathcal B$ dengan $I \cap F$ terbilang, sehingga $F
\setminus P$ termuat di gabungan terbilang banyak himpunan
terbilang itu: jadi terbilang. Karena $F$ tak terbilang,
$P \neq \varnothing$; lalu $P$ tertutup (sebab jika $x \notin P$, maka
maka selang saksi $I$ tak memuat satu titik kondensasi pun:
$I \cap F$ terbilang) dan $P \subseteq F$ (sebab $F$ tertutup:
titik kondensasi khususnya melekat padanya). Lalu $P$
\omterm{prop:b3:galois:perfect}{sempurna}: sebab seandainya $I \cap P = \{x\}$ untuk suatu selang $I$, maka
$I \cap F \subseteq \{x\} \cup (F \setminus P)$ akan
terbilang, dan itu bertentangan dengan $x \in P$. Kini konstruksi pohon
pada pertanyaan 17--19 berjalan kata demi kata di dalam $P$: sebab kepingan
$\bar B(y, \rho) \cap P$ bersifat \omterm{def:b3:topology:compact}{kompak} (sebagai himpunan bagian $\R$ yang tertutup dan terbatas),
setiap pusatnya tak terisolasi di $P$, dan
argumennya tak pernah memakai lebih dari itu. Karena itu $\abs P \geq
\abs{\{0,1\}^\N} = \mathfrak c$, sedangkan $\abs F \leq \mathfrak
c$ secara trivial: jadi $F \subseteq \R$ yang tertutup dan tak terbilang
berkardinalitas tepat $\mathfrak c$. Himpunan tertutup tak dapat
menjadi saksi kegagalan hipotesis kontinum.

\textbf{23.} Misalkan $x \neq y$ di $C$ lalu pilihlah $m \geq 0$ dengan
$3^{-(m+1)} \leq \abs{x - y} < 3^{-m}$. Menurut pertanyaan 1,
angka $b_n(x) = b_n(y)$ bersesuaian untuk $n \leq m$; sedangkan koordinat pertama
$\Phi$ memakai $b_1, b_3, \dots$ dan koordinat keduanya $b_2,
b_4, \dots$, sehingga pada setiap koordinat kedua titik itu berbagi sekurang-kurangnya
$\lfloor m/2\rfloor$ angka biner terdepan:
\[
\norm{\Phi(x) - \Phi(y)}_\infty
\leq \sum_{k > \lfloor m/2\rfloor}2^{-k}
= 2^{-\lfloor m/2\rfloor} \leq \sqrt2\cdot2^{-m/2}
= \sqrt2\,\bigl(3^{-m}\bigr)^\alpha
\leq \sqrt2\,\bigl(3\abs{x-y}\bigr)^\alpha,
\]
dengan memakai $2^{-m/2} = 3^{-m\alpha}$ (ambil logaritmanya); dan
$\sqrt2\cdot3^\alpha = \sqrt2\cdot\sqrt2 = 2$: jadi batasnya
$2\abs{x-y}^\alpha$ pada $C$. Untuk celah $(u, v)$ yang sama: $f$ bersifat afin
di sana, sehingga untuk $t, s \in [u, v]$,
\[
\norm{f(t) - f(s)} = \frac{\abs{t-s}}{v-u}\,
\norm{\Phi(v) - \Phi(u)}
\leq \frac{\abs{t-s}}{v-u}\,2(v-u)^\alpha
\leq 2\abs{t-s}^\alpha,
\]
karena $\frac{\abs{t-s}}{v-u} \leq \bigl(\frac{\abs{t-s}}
{v-u}\bigr)^\alpha$ bila $\abs{t-s} \leq v - u$. Untuk $t <
s$ yang umum: jika $[t, s]$ memotong $C$, misalkan $u'$ dan $v'$ titik terkecil
dan terbesar dari $C \cap [t, s]$ yang \omterm{def:b3:topology:compact}{kompak}; maka $t$
terletak pada sebuah celah (atau pada sebuah titik $C$) yang ujung kanannya
$u'$, sedangkan $s$ pada celah yang ujung kirinya $v'$, dan
\[
\norm{f(t) - f(s)} \leq 2\abs{t - u'}^\alpha
+ 2\abs{v' - u'}^\alpha + 2\abs{s - v'}^\alpha
\leq 6\abs{t - s}^\alpha ;
\]
sedangkan jika $[t, s]$ meleset dari $C$, kasus celah yang sama berlaku. Jadi $f$
bersifat Hölder-$\alpha$ dengan $\alpha = \frac{\ln2}{2\ln3}$.

\textbf{24.} Andaikan $\norm{F(t) - F(s)} \leq
C\abs{t-s}^\alpha$ dengan $F$ surjektif dan $\alpha >
\frac12$. Potonglah $[0,1]$ menjadi $N$ selang $I_1, \dots, I_N$
berpanjang $\frac1N$: maka setiap peta $F(I_j)$ berdiameter $\leq
CN^{-\alpha}$. Dua titik berbeda pada kisi $G_M =
\bigl\{\bigl(\frac iM, \frac jM\bigr)\bigr\}_{0 \leq i, j
\leq M}$ berjarak $\geq \frac1M$, sehingga himpunan berdiameter
$< \frac1M$ memuat paling banyak satu di antaranya. Pilihlah $M =
\lfloor N^\alpha/(2C)\rfloor$ (dengan $N$ besar): maka
$CN^{-\alpha} \leq \frac1{2M} < \frac1M$, dan kesurjektifannya
menempatkan setiap dari $(M+1)^2$ titik kisi itu di suatu
$F(I_j)$, sehingga
\[
\frac{N^{2\alpha}}{4C^2} \leq (M + 1)^2 \leq N,
\]
dan $N^{2\alpha - 1} \leq 4C^2$ gagal untuk $N$ besar bila
$2\alpha > 1$: kontradiksi. Kurva kita, dengan $\alpha \approx
0.3155$, duduk di bawah tembok itu, seperti seharusnya; sedangkan kurva Hilbert
(yang diterima tanpa bukti) bersifat Hölder-$\frac12$, sehingga eksponen kritis
$\frac12$ memang tercapai --- jadi kemulusan Hölder, tak seperti
keinjektifan, merupakan soal derajat, dan $\frac12$ tepat
merupakan perbatasan yang dipaksakan dimensi $2$ pada pemetaan dari
dimensi $1$.

\textbf{25.} Gugusan ke-$m$, yakni $\{\frac1m + \frac1n : n >
m^2\}$, terletak di $\bigl(\frac1m, \frac1m + \frac1{m^2}\bigr]
\subseteq \bigl(\frac1m, \frac1{m-1}\bigr)$ untuk $m \geq 2$,
karena $\frac1{m-1} - \frac1m = \frac1{m(m-1)} >
\frac1{m^2}$ (dan gugusan pertamanya terletak di $(1, \frac32]$):
jadi gugusannya menempati selang yang saling lepas. Titik limit $K$:
di dalam selang ke-$m$ hanya $\frac1m$ (sebab gugusannya
konvergen ke sana dan diskret dalam dirinya sendiri); dan secara global, $0$
(sebab setiap \omterm{def:b3:topology:topology}{persekitaran} $0$ memuat gugusan utuh). Jadi $K'
= \{0\} \cup \{\frac1m\}$, lalu $K'' = \{0\}$ (sebab setiap $\frac1m$
terisolasi di $K'$), dan $K''' = \varnothing$; sedangkan $K$ memuat titik
limitnya, sehingga tertutup, terbatas, terbilang, \omterm{def:b3:topology:compact}{kompak},
dan titik terisolasinya --- yaitu titik gugusan $\frac1m +
\frac1n$ --- \omterm{def:b3:topology:interior}{padat} di $K$: sebab setiap unsur $K'$ merupakan
limitnya, seperti diramalkan mekanisme pertanyaan 21. Induksinya:
$K_1 = \{0\}\cup\{\frac1n\}$ mempunyai $K_1^{(1)} = \{0\}$; lalu diberikan
\omterm{def:b3:topology:compact}{kompak} terbilang $K_j \subseteq [0, 1]$ dengan $K_j^{(j)} =
\{0\}$ dan $K_j^{(j+1)} = \varnothing$, tetapkan
\[
K_{j+1} = \{0\} \cup \bigcup_{m \geq 1}
\Bigl(\frac1m + \varepsilon_m K_j\Bigr),
\]
dengan $\varepsilon_m$ cukup kecil sehingga salinan ke-$m$ terletak
di $\bigl(\frac1m, \frac1{m-1}\bigr)$. Penurunannya bekerja salinan
demi salinan (sebab salinannya hidup di selang terbuka yang saling lepas), sehingga
$K_{j+1}^{(i)} = \{0\} \cup \bigcup_m\bigl(\frac1m +
\varepsilon_mK_j^{(i)}\bigr)$ untuk $i \leq j$; dan pada $i = j$ hal ini
terbaca $\{0\} \cup \{\frac1m\}$, lalu $K_{j+1}^{(j+1)} =
\{0\}$ dan $K_{j+1}^{(j+2)} = \varnothing$. Jadi setiap rank berhingga
muncul. Himpunan \omterm{prop:b3:galois:perfect}{sempurna} merupakan ujung yang lain: $P' = P$, sehingga
penurunannya tak pernah bergerak --- dan Cantor--Bendixson (pertanyaan
22) mengatakan persis bahwa setiap himpunan tertutup terpecah menjadi
inti \omterm{prop:b3:galois:perfect}{sempurna} yang tak terlihat oleh penurunan, dan sisa
terbilang yang dimakan habis oleh penurunan.

\end{solution}
